\chapter{Tunnelling and Underground Works Risk Assessment}
\label{chap:Tunnelling and Underground Works Risk Assessment}

%---------------------------- SECTION ----------------------------
\section{Why this assessment matters in construction}

\paragraph{Tunnelling and underground works represent the most technically demanding and potentially lethal category of civil engineering activity in the UK construction sector. The ground itself is both the medium of construction and the primary source of hazard: it can collapse without warning, flood instantaneously, contain toxic or explosive gases, and resist all attempts at rescue once it has engulfed a worker. The combination of confined space conditions, potential for sudden ground movement, presence of groundwater under pressure, generation of hazardous dusts and fumes in restricted ventilation conditions, and the concentration of heavy plant and services in a geometrically constrained underground environment creates a risk profile that demands the highest levels of specialist competence, engineering control, and regulatory compliance. The Mines Regulations 2014 and the Work in Compressed Air Regulations 1996, together with the broader framework of CDM 2015 and the Confined Spaces Regulations 1997, create specific legal duties for tunnelling projects that go significantly beyond the requirements applicable to surface construction.}

\paragraph{The UK has a substantial ongoing programme of tunnelling work: sewage tunnel rehabilitation, water main replacement, major infrastructure tunnels for rail and road, service tunnels for utilities, and shallow mined works for basement construction in urban environments. Each of these project types presents a characteristic combination of hazards. Deep sewage tunnels intersect the groundwater table and present risks of sudden inundation and noxious gas ingress from the surrounding ground. Urban utility tunnels operate in close proximity to existing infrastructure including gas mains, electricity cables, and existing foundations, creating risks of unplanned service strikes and structural damage to adjacent property. Rail tunnels involve the additional hazard of contact with live traction systems, the specific requirements of Network Rail or the relevant infrastructure manager, and the operational constraints of working in a live railway environment. The risk assessment for any tunnelling project must be specific to its ground conditions, geometry, groundwater regime, and the nature of the underground environment, and must be developed by persons who are genuinely competent in tunnel design and construction.}

\paragraph{Ground conditions in tunnelling are inherently uncertain, and this uncertainty is a defining characteristic of the risk profile. Even with a comprehensive ground investigation programme --- which itself is rarely as comprehensive as the project engineer would wish --- the ground encountered during tunnel construction frequently differs from the conditions assumed in the design. The New Austrian Tunnelling Method (NATM) and observational methods of tunnel construction are specifically designed to manage this uncertainty by monitoring ground behaviour in real time and adjusting the support system accordingly. But the observational method requires a sophisticated monitoring programme, a competent geotechnical team capable of interpreting the data, and a management culture that treats the monitoring data as a trigger for real action rather than as a compliance exercise. The collapse of the Heathrow NATM tunnel under construction in 1994 --- which caused three fatalities and the collapse of substantial surface infrastructure --- remains the defining case study for the consequences of inadequate observational method implementation in UK tunnelling, and its lessons must inform every tunnelling risk assessment prepared today.}

\paragraph{The use of compressed air to stabilise the tunnel face in saturated ground adds a further layer of specific risk that is unique to tunnelling. Compressed air working at pressures above 0.15 bar is subject to the Work in Compressed Air Regulations 1996, which impose detailed requirements for decompression, medical surveillance, emergency procedures, and the provision of a medical lock on site. Decompression sickness (the bends) and barotrauma are occupational health consequences that can cause permanent disability or death if decompression procedures are not strictly followed. The work in compressed air environment requires a level of health and safety management sophistication that is among the most demanding in the entire construction industry.}

\barquote{``The ground does not give warnings that human beings are reliably capable of interpreting. In tunnelling, the consequence of a missed warning is not a near-miss or a minor injury --- it is a collapse that may kill everyone in the heading and prevent their rescue. The monitoring programme is not an administrative requirement: it is the instrument by which the ground communicates its condition to the engineer.''}

\example{NATM Tunnel Face Collapse}{A contractor is constructing a new passenger railway tunnel using the New Austrian Tunnelling Method in London clay at a depth of approximately 20 metres. The ground investigation had characterised the London Clay as stiff and overconsolidated with a design assumption of stand-up time sufficient for the excavation and initial lining sequence. Ground monitoring using settlement points and convergence bolts in the completed tunnel is specified in the design. During a period of continuous excavation, the site team is under programme pressure to advance the heading rapidly. Settlement monitoring results showing progressive settlement at the surface are reported to the site team; the reports are received but not escalated. On a Saturday morning with reduced supervisory presence, the tunnel face collapses over a 6-metre length, and the crown of the completed tunnel section immediately behind the face drops by over 800 mm, causing the collapse of a shallow service tunnel overhead and cracking a gas main running above. No tunnel workers are in the heading at the time of collapse; the Resident Geotechnical Engineer arrives on site and prevents re-entry. The incident results in Prohibition Notices, suspension of all tunnelling works, and a substantial investigation. Investigation concludes that: the monitoring data showing progressive settlement was not reviewed at the frequency required by the specification; the settlement rates that triggered mandatory review thresholds had been reached but not acted upon; and the programme pressure had created a cultural reluctance to raise concerns about ground behaviour. The Heathrow NATM collapse precedent, and the specific warnings in the HSE guidance on NATM tunnelling issued after that incident, had not been reviewed by the project safety team.}

%---------------------------- SECTION ----------------------------
\section{Legal and professional framework}

\paragraph{The legal framework for tunnelling combines general construction health and safety legislation with specific regulations applicable to underground work and compressed air environments. Key frameworks relevant to this chapter include: Health and Safety at Work etc.\ Act 1974; Management of Health and Safety at Work Regulations 1999; Construction (Design and Management) Regulations 2015; Confined Spaces Regulations 1997; Work in Compressed Air Regulations 1996; Mines Regulations 2014 (for qualifying mine-like underground works); Provision and Use of Work Equipment Regulations 1998; Control of Substances Hazardous to Health Regulations 2002; RIDDOR 2013; HSE guidance document HSG185 (Health and Safety in Excavations); and the specific HSE guidance on NATM tunnelling issued by the HSE following the 1994 Heathrow collapse.}

\begin{itemize}
  \bitem{Construction (Design and Management) Regulations 2015}{Impose comprehensive pre-construction planning duties for tunnelling projects, including the requirement for a principal designer to review and coordinate safety in design, a pre-construction information package that must include the ground investigation data and geological interpretive report, a Construction Phase Plan that specifically addresses tunnelling hazards, and a Health and Safety File to be maintained for the completed tunnel structure. Tunnelling projects above the CDM notification threshold (which virtually all major tunnel projects exceed) must be notified to the HSE via the F10 portal.}
  \bitem{Confined Spaces Regulations 1997}{Apply to all underground tunnel environments, which are confined spaces by definition. The CSR 1997 require that the risk of confined space entry is assessed, that a safe system of work is implemented for all entry to the tunnel, and that emergency rescue arrangements are in place before any person enters. The specific application of CSR 1997 to tunnelling is discussed in HSE guidance and in L101 (Safe Work in Confined Spaces).}
  \bitem{Work in Compressed Air Regulations 1996}{Apply whenever work is carried out in compressed air at a gauge pressure above 0.15 bar. The regulations impose requirements for: a Compressed Air Contractor notification to the HSE before work commences; a Compressed Air Safety Supervisor and Compressed Air Engineer to be appointed; a medical lock on site capable of providing emergency recompression treatment; pre-employment and in-employment medical examinations for all compressed air workers; strict decompression tables or computer-controlled decompression; and an emergency response plan for decompression sickness. Employers must not allow any person to work in compressed air who has not passed the medical examination.}
  \bitem{Mines Regulations 2014}{May apply to tunnelling operations that meet the definition of a mine under the Mines Act 1954 --- broadly, excavations below ground for extracting mineral or tunnelling for purposes other than construction of a permanent underground structure. The regulations are primarily relevant to mining operations but the HSE's enforcement approach to certain tunnelling operations with mine-like characteristics applies the Mines Regulations framework.}
  \bitem{COSHH Regulations 2002}{Apply to silica dust generated by tunnel boring in rock, to diesel engine emissions in the tunnel environment, and to any chemical grouts or resins used in ground treatment or tunnel lining operations. The COSHH assessment for tunnelling must address both chemical and physical agents, and the feasibility of engineering controls for dust and fume management in the confined underground environment.}
  \bitem{HSE guidance on NATM tunnelling}{Following the 1994 Heathrow collapse, the HSE published specific guidance on the safety management of NATM tunnelling, with particular emphasis on: the establishment and enforcement of trigger values for monitoring data; the mandatory escalation procedure when trigger values are reached; the maintenance of supervisory oversight during all stages of the excavation and lining cycle; and the prohibition on advancing the heading without compliance with the approved observation cycle. This guidance, while not a statutory instrument, represents the HSE's view of the minimum required standard for NATM projects and will be a central reference in any enforcement investigation following a NATM tunnel incident.}
\end{itemize}

%---------------------------- SECTION ----------------------------
\section{Conducting the assessment using the five-step method}

\subsection{Step 1 --- Identify Hazards}
\paragraph{Tunnelling hazards fall into six primary categories: ground instability and collapse; groundwater inundation; atmospheric hazards (toxic gas, oxygen deficiency, dust, fumes); plant and equipment in confined space; compressed air environment hazards; and surface effects of underground works. Ground instability hazards include: face collapse due to inadequate ground support; crown collapse in completed tunnel sections due to inadequate primary lining; settlement and collapse of surface structures due to ground loss; and piping or heave failure at the invert due to groundwater pressure. Groundwater inundation hazards include: sudden inrush through a high-permeability layer; standpipe failure during dewatering; and progressive seepage undermining primary lining stability. Atmospheric hazards include: methane from organic deposits; carbon dioxide from carboniferous ground; hydrogen sulphide from sulphate-reducing bacteria in wet environments; diesel engine exhaust fumes from plant operating underground; silica and other crystalline dust from rock boring; and cement and grout spray from lining operations. Plant hazards in the confined underground environment include: collision with tunnel boring machine (TBM) components; conveyor belt entanglement; locomotive and tipping car impact; high-voltage electrical shock from TBM power systems; and hydraulic fluid spray from TBM jacking systems.}

\subsection{Step 2 --- Decide Who or What Is Affected}
\paragraph{Persons at risk in tunnelling operations include tunnel workers (known as ``miners'') operating at the tunnel face; operatives managing the muck removal and lining erection train behind the face; TBM maintenance personnel working on or adjacent to the rotating cutterhead; compressed air workers in pressurised environments; surface personnel operating the spoil disposal and material supply systems; and third parties whose surface property may be affected by settlement from the tunnelling operation. Emergency services personnel attending a tunnelling emergency are at particular risk from secondary collapse and from atmospheric hazards in the tunnel environment.}

\subsection{Step 3 --- Evaluate Likelihood and Consequence}
\paragraph{The consequence of a tunnel collapse or inundation involving workers at the heading is catastrophic and virtually always fatal; effective rescue from a major tunnel collapse is, historically, rarely achieved. Ground conditions that deviate from the design assumptions --- softer zones, higher groundwater pressures, or voids from unrecorded former works --- can trigger sudden collapse without prior warning, and the time available for evacuation from the heading to the portal may be seconds. The inherent risk for ground collapse at the face must be rated at the maximum level before engineering controls are applied. Atmospheric hazards from methane or hydrogen sulphide inrush into a tunnel environment can cause multiple fatalities within minutes. The evaluation must be based on the actual ground conditions encountered and the monitoring data, not on the design assumptions.}

\subsection{Step 4 --- Select and Implement Controls}
\paragraph{Controls for tunnelling hazards must be designed and implemented by competent tunnel engineers with specific experience in the ground conditions and tunnel type involved. Ground support must be designed by a geotechnical engineer for the specific ground conditions, and the design must be updated as the actual ground conditions encountered during excavation are logged. The observational method requires a formal monitoring programme with instrument types, installation positions, baseline readings, threshold values, and mandatory response actions defined before excavation commences. Pre-excavation ground treatment --- jet grouting, compensation grouting, dewatering, or compressed air pressure maintenance --- may be required in difficult ground. Atmospheric monitoring must be continuous in tunnel environments where methane or hydrogen sulphide inrush is a credible risk, with fixed-point detectors and alarm systems supplemented by personal monitors carried by all workers in the heading.}

\subsection{Step 5 --- Record and Review}
\paragraph{The tunnel risk assessment is a dynamic document that must be reviewed and updated continuously as the tunnel advances and actual ground conditions are logged. The geological mapping log, monitoring data, and ground support records form a continuous technical record of the tunnel that must be maintained to the highest standard. Any deviation of actual conditions from design assumptions must trigger an immediate review by the geotechnical engineer. The compressed air medical records required by the Work in Compressed Air Regulations 1996 must be maintained for the prescribed period. RIDDOR reporting requirements for tunnelling incidents, including dangerous occurrences underground, must be applied promptly.}

\example{Five-Step Application}{A water company contractor is constructing a 1.5-metre diameter pipeline tunnel in London Clay at 15 metres depth to replace a Victorian water main beneath a residential street. \textbf{Step~1 (Hazards):} Face instability in loose made-ground encountered at the western portal; groundwater inrush risk at the known flint gravel horizon at 12 m depth; diesel exhaust fumes from compressed air loco; silica dust from chalk flints in the London Clay matrix; oxygen deficiency risk in closed headings. \textbf{Step~2 (Who Affected):} Tunnel miners at face, loco drivers, portal crew, surface residents (settlement), utility workers in adjacent services. \textbf{Step~3 (Evaluation):} Face collapse rated Catastrophic/Medium (inherent); groundwater inrush rated Catastrophic/Low (with dewatering); atmospheric hazards rated High/Medium. \textbf{Step~4 (Controls):} Pre-excavation dewatering of gravel horizon; compressed air at 0.5 bar maintained through gravel zone (Work in Compressed Air Regulations 1996 compliance confirmed); continuous four-gas monitoring in heading; forced ventilation with dilution calculation for diesel; clay face support using steel arches at 0.75 m centres; surface settlement monitoring at 2 m grid with 10 mm amber trigger and 20 mm red trigger; emergency procedure tested before heading commenced. \textbf{Step~5 (Record):} Geological mapping log updated daily; monitoring data reviewed by geotechnical engineer weekly and following any amber trigger; compressed air medical records maintained; daily confined space permit for heading access.}

%---------------------------- SECTION ----------------------------
\section{Applying the hierarchy of controls}

\paragraph{The hierarchy of controls in tunnelling begins with ground engineering rather than personal protection, because in the underground environment no level of PPE can protect against the energy release of a tunnel collapse or an inrush of water.}

\begin{itemize}
  \bitem{Elimination}{Design the tunnel alignment to avoid the highest-risk ground conditions where practicable: avoid zones of known natural or manmade cavities, avoid the wettest groundwater horizons if the alignment can be adjusted, and avoid operating within close proximity of existing live services or foundations where an alternative route is available. In urban environments, route selection is constrained, but even marginal improvements in alignment can significantly reduce the ground risk.}
  \bitem{Substitution}{Substitute higher-risk excavation methods with lower-risk alternatives where ground conditions permit. Closed-face tunnel boring machines (TBMs) with earth pressure balance or slurry pressure balance systems substitute the open-face methods that require more intensive manual ground support. Microtunnelling eliminates manned entry to the face entirely. Pre-excavation ground treatment (jet grouting, compensation grouting, ground freezing) substitutes weak ground with treated ground of predictable engineering properties before excavation begins.}
  \bitem{Engineering controls}{Ground support system designed for the actual ground conditions encountered; monitoring programme with defined trigger values and mandatory escalation procedure; pre-excavation dewatering or compressed air to manage groundwater; continuous atmospheric monitoring with fixed-point detectors and personal monitors; forced ventilation with calculated dilution factor for diesel fumes; emergency communication system from face to surface; and escape capsule or escape breathing apparatus at the face for use in the event of atmospheric emergency.}
  \bitem{Administrative controls}{Confined space entry permit for all access to the tunnel; daily briefing for tunnel workers on ground conditions, monitoring status, and emergency procedures; geological mapping log maintained by competent geotechnical engineer; monitoring data reviewed at defined frequency with mandatory escalation when trigger values are reached; strict application of the Work in Compressed Air Regulations 1996 decompression procedures; no advancement of heading without compliance with the approved support and observation cycle.}
  \bitem{PPE}{Hard hat, hearing protection (in areas above 85 dB(A)), respiratory protection (P3 filter FFP3 disposable half-mask or better for silica dust environments), and personal four-gas monitor carried by all workers in the heading. Compressed air workers must wear the specific PPE specified in the health surveillance and decompression protocols.}
\end{itemize}

%---------------------------- SECTION ----------------------------
\section{Cadence of assessment and review triggers}

\paragraph{The tunnelling risk assessment must be reviewed continuously as the tunnel advances, and formally reviewed at defined milestones and following specific trigger events.}

\begin{itemize}
  \bitem{At the commencement of each new ground type}{When the geological mapping log indicates a transition to a new ground stratum, the geotechnical engineer must review the ground support design and the risk assessment for the new conditions before the heading is advanced. Transitions to weaker, wetter, or more gaseous ground may require significant changes to the support system or excavation method.}
  \bitem{When a monitoring trigger value is reached}{When settlement, convergence, or groundwater monitoring data reaches a defined amber or red trigger value, the mandatory escalation procedure must be activated immediately. Work at the heading must cease; the geotechnical engineer must review the data and the ground conditions; and the risk assessment must be reviewed and, if necessary, revised before work resumes. The trigger value system is only effective if triggers result in action --- a monitoring system that generates data but produces no response to threshold exceedance is worse than no monitoring at all.}
  \bitem{Following any unplanned event underground}{Any unplanned ground movement, water ingress, atmospheric alarm actuation, or plant failure underground must trigger an immediate formal review of the risk assessment, the support system design, and the excavation method. The tunnel heading must be secured before the review is completed.}
  \bitem{Following an incident or near miss}{Any injury, near miss, or dangerous occurrence in the tunnel must trigger a full formal investigation and risk assessment review before work resumes. RIDDOR reporting obligations must be assessed immediately.}
  \bitem{Monthly formal review of tunnel risk assessment}{The tunnel risk assessment as a whole must be formally reviewed monthly by the tunnel safety manager, the geotechnical engineer, and the principal contractor's project manager, taking account of the geological log to date, the monitoring data trend, any incidents or near misses, and any changes to the planned programme.}
\end{itemize}

%---------------------------- SECTION ----------------------------
\section{Common failure modes}

\begin{itemize}
  \bitem{Monitoring data collected but not acted upon}{The classic failure mode identified in NATM tunnel incidents is the generation of settlement and convergence monitoring data showing progressive ground movement, which is not reviewed at the frequency or with the urgency required, and does not trigger the mandatory response actions specified. The data exists; the threshold has been reached; no action follows. This failure mode is not a technical failure --- it is a management and cultural failure, and it is the failure mode that killed workers in the Heathrow NATM incident.}
  \bitem{Ground conditions deviate from design but support system is not revised}{The geological mapping log records a transition to weaker ground, but the ground support system designed for the original design conditions is not updated by the geotechnical engineer. The heading advances into weaker ground with inadequate support, and collapse follows. The requirement for the observational method to actually result in revised support designs when ground conditions deviate is fundamental to its safety function.}
  \bitem{Compressed air decompression procedures not followed under programme pressure}{Programme pressure and worker incentive structures create a cultural pressure to reduce decompression time. Even small reductions in decompression duration increase the risk of decompression sickness. The Work in Compressed Air Regulations 1996 decompression tables represent the minimum safe decompression time for the working pressure and duration; any reduction below the table values is a direct legal breach and a physiological risk.}
  \bitem{Atmospheric monitoring in the heading inadequate for the gas environment}{Atmospheric monitoring that detects oxygen deficiency but does not measure methane, or that measures methane but does not have a sensor positioned at the highest point of the heading where methane accumulates, does not provide the protection intended. The specific gases that may be present in the ground conditions being excavated must be assessed and appropriate sensors provided.}
  \bitem{Emergency rescue plan not pre-planned and practised}{Emergency rescue from a tunnel collapse is technically difficult and time-critical. A rescue plan that exists only on paper, has never been rehearsed, and has not been shared with the local emergency services is not a rescue plan in any operational sense. Tunnelling projects must conduct regular emergency exercises, brief the local fire service on the underground layout, and maintain rescue equipment at the portal in operational condition.}
\end{itemize}

\example{Failure Pattern}{A contractor is advancing a sewer tunnel in weathered Mercia Mudstone using drill-and-blast in a dry section, transitioning to a wet zone identified in the ground investigation. The geotechnical engineer notes in the geological log that the mudstone ahead is fractured and shows signs of water pressure behind the face. The site team, under programme pressure, continues without pre-excavation dewatering and without increasing the steel rib support spacing. After two further advance rounds, a zone of fractured mudstone with water pressure behind it fails suddenly, flooding the heading. Three workers are in the heading at the time. Rescue operations are hampered by the inrush of water and the confined access. One worker escapes; two workers are trapped and, after a prolonged rescue operation, one survives with injuries and one is recovered fatally. Investigation confirms that the geological log entries identifying the risk of inundation were present in the record; the dewatering specification was not implemented; and the geotechnical engineer's field notes, recommending pre-drainage before advancing, had not been escalated to the project manager for a decision.}

%---------------------------- CHAPTER SUMMARY ----------------------------
\newpage
\section{Chapter Summary}

\begin{table}[H]
\centering
\begin{tabularx}{\textwidth}{>{\raggedright\arraybackslash}p{3.2cm}X}
\toprule
\textbf{Assessment Element} & \textbf{Operational Requirement} \\
\midrule
Legal/Professional Basis & CDM 2015; Confined Spaces Regulations 1997; Work in Compressed Air Regulations 1996; Mines Regulations 2014; COSHH 2002; PUWER 1998; HSE NATM guidance; Health and Safety at Work etc.\ Act 1974. \\
Method & Apply the full five-step process and document inherent/residual risk. Engage a competent tunnel engineer. Apply the observational method with defined monitoring triggers and mandatory escalation procedures. \\
Controls & Ground support designed for actual conditions and updated as ground type changes; monitoring programme with trigger values and mandatory response; compressed air compliance with Work in Compressed Air Regulations 1996; continuous atmospheric monitoring with personal monitors; emergency rescue plan rehearsed and shared with emergency services. \\
Cadence & On transition to new ground type; when monitoring trigger reached; following any unplanned underground event; following any incident; monthly formal review. \\
Assurance & Geological mapping log; monitoring data records and review notes; CDM F10 notification; Work in Compressed Air medical records; RIDDOR reports; emergency exercise records. \\
\bottomrule
\end{tabularx}
\end{table}

\begin{itemize}
  \bitem{Key takeaway 1}{Tunnelling risk assessment must be developed and maintained by a competent tunnel engineer with specific experience in the relevant ground conditions and tunnel type. Generic construction risk assessment competence is not sufficient for tunnelling.}
  \bitem{Key takeaway 2}{The observational method only delivers its safety function if monitoring data is reviewed at the required frequency, trigger values result in mandatory escalation, and the geotechnical engineer has the authority and the management support to stop work when threshold values are reached.}
  \bitem{Key takeaway 3}{Work in compressed air above 0.15 bar gauge is subject to the Work in Compressed Air Regulations 1996, which impose specific requirements for decompression, medical surveillance, emergency recompression facilities, and notification to the HSE. Decompression tables must not be shortened under programme pressure.}
  \bitem{Key takeaway 4}{Atmospheric monitoring in the heading must be appropriate for the specific gas hazards presented by the ground conditions being excavated. Fixed-point detectors must be positioned to detect the lightest gases (methane) at the crown and the heaviest (CO\textsubscript{2}, H\textsubscript{2}S) at the invert.}
  \bitem{Key takeaway 5}{Emergency rescue from a tunnel collapse or inrush is technically difficult and time-limited. The rescue plan must be pre-planned, rehearsed, communicated to the emergency services, and supported by rescue equipment maintained in operational condition at the tunnel portal.}
\end{itemize}

\barquote{In Chapter~\ref{chap:Roofing, Cladding and Facade Risk Assessment}, we examine the risk assessment requirements for roofing, cladding, and façade works --- operations that combine working at height with specialist materials handling, weather exposure, and the interaction with structural and waterproofing systems at the building envelope.}
